\frontsection{Подготовка среды и исходных материалов}
\ifdefempty{\CompanionURL}{}{%
\noindent Электронное приложение к практикуму:\par
\noindent\expandafter\url\expandafter{\CompanionURL}\par\medskip}
\subsection*{Структура файлов}
Держите исходные данные разных работ в отдельных каталогах. Файлы с одинаковыми именами относятся к разным задачам и не взаимозаменяемы.
\begin{lstlisting}[language={}]
project/
  code/                         # common teaching modules
  labs/
    01_classification/
      lab1_activity.ipynb
      data/                     # train.csv, test.csv
    02_regression/
      lab2_housing.ipynb
      data/                     # CSV and data_description.txt
    03_cnn/
      lab3_flowers.ipynb
      data/                     # five class folders
\end{lstlisting}

\noindent\begin{tabularx}{\textwidth}{@{}p{0.22\textwidth}Y@{}}
\toprule
Работа & Что должно быть в данных \\
\midrule
№\,1: активность & В обеих таблицах: числовые признаки, метка \file{Activity}, идентификатор участника (в коде --- \file{subject}). \\
№\,2: жильё & В обучающей таблице: \file{Id}, признаки и \file{SalePrice}; в таблице прогнозирования цены нет. Описание признаков находится в \file{data\_description.txt}. \\
№\,3: цветы & Пять каталогов видов с изображениями. Индексы классов определяются фактическим отображением \file{class\_to\_idx}. \\
\bottomrule
\end{tabularx}
\par\medskip
Преподаватель размещает выборки в папках работ или выдаёт их отдельно. Если в \file{data} пока только инструкция, добавьте необходимые файлы. Набор изображений указан в исходном задании: \href{https://www.kaggle.com/datasets/alxmamaev/flowers-recognition}{Flowers Recognition}. Количество строк, признаков и изображений определяйте по своей версии файлов.

\subsection*{Что такое ноутбук и как его запускать}
Ноутбук Jupyter (файл \file{.ipynb}) состоит из текстовых ячеек с пояснениями и ячеек Python. Выполняйте их сверху вниз. Переменные предыдущих ячеек остаются в памяти: если выполнить только нижнюю ячейку после изменения данных, можно случайно использовать старый результат. Перед итоговой проверкой перезапустите среду исполнения (kernel/runtime) и выполните все ячейки заново.

GitHub показывает содержимое ноутбука и хранит файлы; код исполняется в Jupyter Notebook или JupyterLab на компьютере студента либо учебном сервере. Для этого комплекта нужен весь проект с папкой \file{code}, а не только один \file{.ipynb}. В \file{START\_HERE.md} электронного приложения приведён порядок первого запуска. На GitHub сохраняйте себе копию материалов согласованной версии.

\subsection*{Программная среда}
Для работ №\,1 и №\,2 нужны Python, NumPy, pandas, scikit-learn и Matplotlib; для второго примера также TensorFlow/Keras. Для работы №\,3 нужны PyTorch, совместимая с ним torchvision и Pillow. Инструкции по проверенному локальному окружению находятся в \file{code/README.md}.

CPU --- центральный процессор компьютера. GPU --- ускоритель, который может быстрее выполнять многие операции нейросетей; его наличие не меняет смысла задания. В готовом учебном окружении сначала проверьте версии библиотек. При самостоятельной установке согласуйте версии Python, PyTorch, torchvision и доступного ускорителя по документации. Для работы на GPU нужна подходящая сборка библиотек.

\begin{lstlisting}
import sys
from importlib.metadata import version
print(sys.version)
for name in ("numpy", "pandas", "scikit-learn"):
    print(name, version(name))
\end{lstlisting}

В третьей работе ускоритель существенно сокращает время обучения. Проверка \file{torch.cuda.is\_available()} сообщает о доступности CUDA, а не о достаточном объёме памяти. Короткую проверку небольшого варианта можно выполнить на CPU; для сравнения тяжёлой исходной сети заранее оцените ресурсы.

\subsection*{Пути в ноутбуке}
Первая ячейка кода находит корень проекта по папкам \file{code} и \file{labs}, затем выбирает каталог своей работы. Поэтому путь \file{data/train.csv} ведёт к данным именно этой работы, а \file{runs} --- к её результатам. Запускать ноутбук можно из корня проекта или каталога работы. Сохраняйте структуру скачанного комплекта; для нового опыта выбирайте новое имя папки результатов.

\begin{keybox}{Перед запуском}
Проверьте, что открываете данные нужной работы; можете прочитать таблицу или изображение; понимаете смысл целевого ответа; знаете, какие файлы будут сохранены. Зафиксируйте версии библиотек и seed. Не сравнивайте результаты, полученные на незаметно разных разбиениях.
\end{keybox}
